% English translation of prelude.tex, lines 154--183.
% Source: Methods of Algebra, Volume 2, commit
% 9a5803ff77dd3257484cb177f851a73770a59dd3 (CC BY 4.0).
% stable-unit-id: o014.aljabr2.prelude.structure-overview

% segment-id: o014.li2.prelude.structure-overview.h001
\section*{Overview of the Structure}
\addcontentsline{toc}{section}{Overview of the Structure}
\hypertarget{o014.aljabr2.prelude.structure-overview}{ }

% segment-id: o014.li2.prelude.structure-overview.p001
Readers are advised to begin with the Conventions at the end of the introduction, in order to familiarize themselves with the notation and conventions used throughout the book. Beyond that, beginners are discouraged from reading strictly in sequence unless they are already highly receptive to abstract methods. Each chapter opens with a detailed summary and reading guide.

% segment-id: o014.li2.prelude.structure-overview.p002
First comes the Inner Part, which lays the foundations.

\begin{description}
	% segment-id: o014.li2.prelude.structure-overview.i001
	\item[Chapter 1: Supplements to Category Theory] Reviews and rounds out the requisite category theory, emphasizing strict morphisms, kernels and cokernels in additive categories, Kan extensions, Gabriel--Zisman localization, and related topics.
	% segment-id: o014.li2.prelude.structure-overview.i002
	\item[Chapter 2: Abelian Categories] Develops the basic theory of abelian categories, supplemented by material on Grothendieck categories and other topics.
	% segment-id: o014.li2.prelude.structure-overview.i003
	\item[Chapter 3: Complexes] Discusses basic notions such as complexes, double complexes, homotopies, and mapping cones in additive categories, as well as cohomology in abelian categories. The second half treats derived functors from the classical viewpoint, including special cases such as $\Ext$, $\Tor$, and $\lim^1$. To handle unbounded complexes, it concludes with K-injective and K-projective resolutions.
	% segment-id: o014.li2.prelude.structure-overview.i004
	\item[Chapter 4: Triangulated and Derived Categories] Introduces the general notions of triangulated categories and triangulated functors, then concentrates on derived categories of abelian categories and derived functors between derived categories. The classical $\Ext$ and $\Tor$ are upgraded accordingly to $\RHom$ and $\otimesL$.
	% segment-id: o014.li2.prelude.structure-overview.i005
	\item[Chapter 5: Spectral Sequences] Begins with the general definition of spectral sequences and exact couples, then specializes progressively to the spectral sequences of filtered complexes and double complexes and to their applications in algebra. Exact couples are not absolutely necessary for the spectral sequence of a filtered complex.
\end{description}

% segment-id: o014.li2.prelude.structure-overview.p003
Next comes the Outer Part, devoted to applications or further developments.
\begin{description}
	% segment-id: o014.li2.prelude.structure-overview.i006
	\item[Chapter 6: Group Homology and Cohomology] Studies group homology and cohomology, together with a variety of related computations and examples. It also treats Tate cohomology, the cohomology of profinite groups, and nonabelian cohomology, each of which plays an important role in applications.
	% segment-id: o014.li2.prelude.structure-overview.i007
	\item[Chapter 7: Monad Theory] Starts from the general notion of an ``algebra'' in a monoidal category and first discusses differential graded structures related to complexes. After treating examples that include coalgebras, bialgebras, and Hopf algebras, it turns to Beck's monadicity theorem. The second half considers module theory, introduces Morita theory, and applies monadicity to flat descent for modules and to Galois descent.
	% segment-id: o014.li2.prelude.structure-overview.i008
	\item[Chapter 8: Simplicial Methods] Covers the general theory of simplicial objects and simplicial sets, the Dold--Kan correspondence, and the bar construction formulated in the language of monads. The final sections concern the generalized Eilenberg--Zilber theorem and a topological interpretation of mapping cones.
	% segment-id: o014.li2.prelude.structure-overview.i009
	\item[Chapter 9: Duality] Begins with duality in monoidal categories and, after basic examples, studies reconstruction theorems and Tannakian categories. This chapter is not directly concerned with complexes.
\end{description}

% segment-id: o014.li2.prelude.structure-overview.p004
The original plan called for the Outer Part to include the foundations of representation theory, which seemed likely to play a natural role in the overall structure. Late in the process, the idea was abandoned because the book was already too long and several mature textbooks were available.

% segment-id: o014.li2.prelude.structure-overview.p005
The appendices are as follows.
\begin{description}
	% segment-id: o014.li2.prelude.structure-overview.i010
	\item[Appendix A: Further Material on Abelian Categories] Collects material directly or indirectly related to abelian categories, either cited in the main text or developed as a side topic.
	% segment-id: o014.li2.prelude.structure-overview.i011
	\item[Appendix B: An Introduction to Ind- and Pro-objects] Starts from the theory of profinite groups, naturally introduces ind- and pro-objects, and studies the Ind-construction for categories and functors. As an application, it then proves the Freyd--Mitchell embedding theorem for abelian categories.
\end{description}

% segment-id: o014.li2.prelude.structure-overview.p006
Finally, the author's personal website provides the latest supplementary information, including errata, for this book and the author's other works. Readers are invited to follow the updates there and to write to the author if they find any errors in the book.
